\chapter{Schwer verstandene Phänomene als zusätzliche Prüfsteine}
\label{ch:shadow-open}
Offene Physik ist kein freier Platz, in den jede neue Idee hineinpasst. Sie ist oft durch viele bekannte Eigenschaften eng begrenzt. Der nützliche Schritt lautet deshalb: Welcher Teil des Phänomens ist beobachtet, welcher Mechanismus fehlt, und welche Antwort müsste unser Modell zusätzlich liefern?

\section{Dunkle Materie: unsichtbar heißt nicht frei wählbar}
Gravitationslinsen und die Verteilung sichtbarer Materie in kollidierenden Haufen liefern eine wichtige Beobachtung: Der aus der Gravitation erschlossene Beitrag muss nicht dem heißen sichtbaren Gas folgen. Die Bullet-Cluster-Analyse ist ein prominentes Beispiel. Sie identifiziert damit noch kein bestimmtes mikroskopisches Teilchen. Zusammen mit anderen kosmologischen Daten ist sie eine starke Anforderung an jede Erklärung. \src{F15; E16}

Im Universalraum könnte man an schwach auslesbare Sektoren denken. Das ist zunächst eine Hypothese. Der bloße Umstand, dass eine Komponente in einem Schatten unsichtbar bleibt, gibt ihr keine gravitative Wirkung, passende Dichte oder lange Lebensdauer. Ein konkreter Kandidat müsste gleichzeitig seine geringe elektromagnetische Kopplung, seine Gravitation, seine Entstehung und seine Verteilung erklären. Ein lokales Register, das ein Beobachter gerade nicht ausliest, ist deshalb nicht schon dunkle Materie.

\section{Dunkle Energie: ein kleiner Rest mit großen Folgen}
Die beschleunigte kosmische Expansion verlangt eine konsistente Erklärung im Gesamtmodell. Für TFPT ist besonders schwierig, eine passende effektive Vakuumwirkung zu erhalten, ohne sie bei jedem neuen Beitrag nachzujustieren. Der Mechanismus müsste seine Größe und Stabilität gegenüber Korrekturen erklären.

DESI DR2 untersucht unter anderem eine zeitabhängige Zustandsgleichung. Die Präferenz hängt von den kombinierten Datensätzen und dem gewählten Modell ab; sie ist nicht mit einer gesicherten Entdeckung dynamischer dunkler Energie gleichzusetzen. Der hier herangezogene Bericht ist die veröffentlichte DR2-Auswertung von 2025, im September 2026 geprüft. \src{F16}

Ein \glqq Informationsdruck\grqq{} des Universalraums würde erst dann etwas erklären, wenn ein wirksamer Energie-Impuls-Tensor folgt: mit Erhaltungsgesetz, Zustandsgleichung, stabilen Störungen und überprüfbarer Entwicklung. Ohne diesen Schritt bleibt das Wort ein Bild. Ebenso darf eine Spannung zwischen kosmologischen Messwegen nicht von vornherein als Beweis einer bestimmten neuen Dynamik behandelt werden.

\section{Schwarze Löcher: Information kann die Aufteilung selbst betreffen}
Die Schwarzlochfrage verbindet Quanteninformation, Gravitation und die Wahl zugänglicher Teilsysteme besonders eng. In bestimmten berechenbaren Modellen verändern Quanteneffekte die extremale Fläche und damit die Zuordnung rekonstruierbarer Information. Die verallgemeinerte Entropie enthält einen Flächenbeitrag und einen Quantenfeldbeitrag. Das ist erheblicher theoretischer Fortschritt in diesen Modellen, aber noch keine vollständige Beschreibung eines realistischen verdampfenden astrophysikalischen schwarzen Lochs. \src{F17}

Für die Universalraum-Idee ist daran interessant, dass der \emph{Zugang} zu Information von Geometrie und Dynamik abhängt. Das passt zur operationalen Definition von Schatten. Eine Ableitung müsste jedoch viel konkreter werden: Woher stammt ein Horizont, warum tritt ein Flächenterm auf, und welche Algebra ist jeweils zugänglich? Ein endliches Aufzeichnungsregister mit Entropie genügt nicht zur Herleitung der Bekenstein--Hawking-Beziehung.

\section{Gravitation muss sehr verschiedene Inhalte gleich behandeln}
Äquivalenztests wie MICROSCOPE prüfen die Universalität des freien Falls für unterschiedliche Materialien mit hoher Genauigkeit. Eine neue Quelle für Raumzeit muss daher erklären, warum verschiedene innere Zusammensetzungen auf dieselbe effektive Geometrie reagieren. \src{F18}

Ein gemeinsamer Kausalkegel ist dafür ein Anfang. Wenn aber jeder TFPT-Sektor seine eigenen Ausbreitungsgeschwindigkeiten oder seine eigene gravitative Kopplung behalten würde, wäre die allgemeine Universalität noch nicht gewonnen. Die Aufgabe lautet nicht nur, ein Spin-2-ähnliches Muster zu finden, sondern eine konsistente, universelle Kopplung mit den bekannten Niedrigenergiegrenzen zu erhalten.

\section{Weitere offene Stellen mit unterschiedlicher Beweiskraft}
\begin{longtable}{L{32mm}L{56mm}L{60mm}}
\toprule\textbf{Thema}&\textbf{Was den Ansatz einschränkt}&\textbf{Was eine Erklärung liefern müsste}\\\midrule\endhead
Zeitpfeil&Beständige Aufzeichnungen und irreversible Makroprozesse bei oft reversiblen Grundgleichungen.&Zustandsauswahl, Anfangskorrelationen, Zugriff und Entropiefluss; keine bloße Benennung eines inneren Parameters als Zeit.\\
Messproblem&Interferenz und stabile Messergebnisse müssen zusammen beschrieben werden.&Ein vollständiges Messmodell samt Wahrscheinlichkeiten oder eine abweichende Dynamik mit unterscheidbarer Vorhersage.\\
Massen und Familien&Gemessene Teilchenspektren und Mischungen, einschließlich Neutrino-Oszillationen.&Die Auswahl der Sektoren und ihrer Parameter mit Fehlerrechnung; passende kleine Zahlen allein reichen nicht.\\
Materieüberschuss&Die heutige Materie-Antimaterie-Asymmetrie verlangt eine kosmologische Entstehungsgeschichte.&Geeignete Symmetrieverletzung und Nichtgleichgewichtsdynamik mit quantitativer Ausbeute.\\
Starke CP-Frage&Im starken Sektor muss eine mögliche CP-verletzende Wirkung sehr klein sein.&Einen Schutz- oder Auswahlmechanismus samt Konsequenzen; ein Vorzeichenmuster im internen Gitter ist noch kein solcher Mechanismus.\\
Quantenfelder und Gravitation&Quantenwahrscheinlichkeiten, Lokalität und die erfolgreichen klassischen Näherungen müssen vereinbar bleiben.&Eine gemeinsame Dynamik mit kontrollierten Grenzwerten und experimentell zugänglichen Abweichungen.\\\bottomrule
\end{longtable}
\src{Einordnung; quantitative Referenzdaten in E14--E16, zusätzliche Primärquellen F07--F19; F19: Neutron-Dipolmessung als CP-Prüfstein}

Diese Liste behauptet nicht, alle Probleme seien gleich unverstanden. Quantenmechanik besitzt extrem erfolgreiche Vorhersageregeln, obwohl ihre Interpretation diskutiert wird. Dunkle Materie hat eine starke beobachtete gravitative Rolle, obwohl ihre mikroskopische Natur offen ist. RH ist dagegen ein exakt formuliertes mathematisches Problem. Die verschiedenen Arten von Unwissen dürfen nicht zu einem einzigen scheinbaren Beleg verschmelzen.

\begin{bild}{Die verdichtete Lektion}
Ein gemeinsamer Ursprung wäre besonders überzeugend, wenn ein und dieselbe getroffene Wahl mehrere Aufgaben gleichzeitig löst: etwa lokale Quantenkorrelationen, stabile Ausbreitung und eine bestimmte arithmetische Spur. Ein neues freies Einstellrad pro Phänomen würde den gemeinsamen Ursprung noch nicht erklären.
\end{bild}
